\section{Methodology}
\label{sec:methodology}

\subsection{Experimental Pipeline}

Each experiment follows the two-stage pipeline below.
Stage 1 (channel transmission):
\[
  \text{text} \xrightarrow{\text{UTF-8}} \text{bits}
  \xrightarrow{\text{ECC encode}} \text{codeword}
  \xrightarrow{\mathrm{BSC}(p)} \tilde{\text{codeword}}
\]
Stage 2 (recovery):
\[
  \tilde{\text{codeword}}
  \xrightarrow{\text{ECC decode}} \hat{\text{bits}}
  \xrightarrow{\text{UTF-8}} \hat{\text{text}}
  \xrightarrow{[\text{LLM}]} \hat{\hat{\text{text}}}
\]
The LLM correction stage is optional and applied to the ECC output (or directly
to the raw BSC output in the LLM-only condition).

\subsection{Text Corpus}

We use a fixed corpus of 40 English sentences drawn from technical and
general-language domains, including descriptions of communication theory,
information-theoretic concepts, and general knowledge statements.  Sentences
range from 45 to 110 characters.  The corpus is held fixed across all
experiments to enable direct comparison.

\subsection{Evaluation Metrics}

We evaluate at four levels of abstraction:
\begin{itemize}
  \item \textbf{BER} (Bit Error Rate): fraction of bits differing from original, measured on the recovered bit string before UTF-8 decoding.
  \item \textbf{CER} (Character Error Rate): edit distance between original and recovered text, normalized by original length.
  \item \textbf{WER} (Word Error Rate): edit distance at the word level, normalized by original word count.
  \item \textbf{BLEU}: 4-gram precision with brevity penalty \cite{papineni2002}, measuring $n$-gram overlap.
\end{itemize}
BER captures channel-level performance; CER/WER/BLEU capture human-readable quality.

\subsection{LLM Prompting Strategies}

We evaluate three prompting strategies for the LLM correction stage:
\begin{enumerate}
  \item \emph{Minimal}: a bare instruction to correct the corrupted text.
  \item \emph{Contextual}: includes the known noise level $p$ as context.
  \item \emph{Exemplar}: provides two (corrupted, original) examples before the target.
\end{enumerate}
All experiments use \texttt{claude-sonnet-4-6} with a cached system prompt to
reduce API cost.

\subsection{Shannon Bound Analysis}

For each method and each $p$ value, we compare the achieved information rate
(estimated as $1 - H_b(\text{BER}_\text{post})$) against the Shannon capacity
$C(p)$ and the code rate $R$ of the ECC.  A reliable code must satisfy
$R \leq C(p)$; empirical results above this bound indicate residual errors not
yet corrected, and the rate estimate is a lower bound on achievable performance.

\subsection{Reproducibility}

All experiments use fixed random seeds (base seed 42, per-trial offset by trial
index).  LLM responses are cached in a JSONL file keyed by a SHA-256 hash of
the input text, model identifier, and prompt strategy, ensuring that re-running
experiments does not incur additional API cost.
